% Source-faithful mathematical blueprint of Hartshorne, Algebraic Geometry (1977), Appendix B.
% Authorial exposition, examples, and exercises are omitted; mathematical definitions, statements,
% formulas, and proof arguments are retained.
\section{The Associated Complex Analytic Space}
\label{ha-appB-sec1}
\source{hartshorne-algebraic-geometry:B.1}

\begin{construction}[Analytification]
\source{hartshorne-algebraic-geometry:B.1}
If $X$ is of finite type over $\CC$, cover it by $Y_i=\operatorname{Spec}A_i$
with $A_i=\CC[x_1,\ldots,x_n]/(f_1,\ldots,f_q)$.  Regard the $f_j$ as
holomorphic functions, form the analytic zero spaces $(Y_i)_h\subset\CC^n$,
and glue them using the algebraic transition maps.  The resulting analytic
space is $X_h$, functorially in $X$.  A coherent sheaf analytifies by applying
the same construction to local finite cokernel presentations.  The
construction preserves separatedness/Hausdorffness, connectedness,
reducedness, smoothness/being a complex manifold, and properness.  There are
natural maps
\[
H^i(X,\mathcal F)\longrightarrow H^i(X_h,\mathcal F_h).
\]
\end{construction}

\section{Comparison of the Algebraic and Analytic Categories}
\label{ha-appB-sec2}
\source{hartshorne-algebraic-geometry:B.2}

\begin{theorem}[GAGA]
\label{ha-appB-thm2-1}
\dcref{B.2.1}
If $X$ is projective over $\CC$, analytification gives an equivalence between
coherent algebraic sheaves on $X$ and coherent analytic sheaves on $X_h$;
for every coherent $\mathcal F$ the comparison maps
\[
\alpha_i:H^i(X,\mathcal F)\xrightarrow{\sim}H^i(X_h,\mathcal F_h)
\]
are isomorphisms.
\source{hartshorne-algebraic-geometry:B.2.1}

\end{theorem}
\begin{proof}
Embed $X$ in $\mathbf P^N$.  Analytification of a finite presentation gives a coherent analytic
sheaf, and Serre's finiteness theorems A and B identify its Cech cohomology with the algebraic Cech
complex on the standard affine cover.  The resulting comparison maps are isomorphisms and are
compatible with tensor products and morphisms, yielding the equivalence.
\end{proof}

\begin{theorem}[Chow]
\label{ha-appB-thm2-2}
\dcref{B.2.2}
Every compact analytic subspace of $\mathbb P^n_h$ is the analytification of
a unique closed algebraic subscheme of $\mathbb P^n$.
\source{hartshorne-algebraic-geometry:B.2.2}

\end{theorem}
\begin{proof}
This is stated as a corollary of Serre's theorem in the source; no separate
proof is supplied.
\end{proof}

\section{When Is a Compact Complex Manifold Algebraic?}
\label{ha-appB-sec3}
\source{hartshorne-algebraic-geometry:B.3}

\begin{theorem}[Riemann]
\label{ha-appB-thm3-1}
\dcref{B.3.1}
Every compact complex manifold of dimension one is projective algebraic.
\source{hartshorne-algebraic-geometry:B.3.1}
\end{theorem}

\begin{proof}
The source explains that the first step uses hard analysis (Weyl's or
Gunning's argument) to produce a nonconstant meromorphic function; the second,
elementary Riemann-existence step gives a finite map to $\mathbb P^1$ and
algebraicity.  No full proof is supplied.
\end{proof}


\begin{theorem}[Generalized Riemann existence]
\label{ha-appB-thm3-2}
\dcref{B.3.2}
If $X$ is a normal scheme of finite type over $\CC$, and $X'$ is a normal
complex analytic space with a finite morphism $f:X'\to X_h$, then there is a
unique normal scheme $X'$ and finite morphism $g:X'\to X$ such that
$X'_h\simeq X'$ and $g_h=f$.  Consequently
\[
\pi_1^{\mathrm{alg}}(X)\simeq\widehat{\pi_1(X_h)}.
\]
\source{hartshorne-algebraic-geometry:B.3.2}
\end{theorem}
\begin{proof}
The source attributes this generalized theorem to Grauert--Remmert and
Grothendieck (SGA 1) and does not supply a proof.
\end{proof}

\begin{proposition}[Siegel]
\label{ha-appB-prop3-3}
\dcref{B.3.3}
For a compact complex manifold $X$ of dimension $n$, the meromorphic function
field has transcendence degree at most $n$ over $\CC$; if equality holds, it is
a finitely generated extension of $\CC$.  If $X=X_h$ is algebraic, then this
field is the algebraic rational function field.
\source{hartshorne-algebraic-geometry:B.3.3}

\end{proposition}
\begin{proof}
The source cites Siegel for this proposition and does not give its proof.
\end{proof}

\begin{theorem}[Chow--Kodaira]
\label{ha-appB-thm3-4}
\dcref{B.3.4}
Every compact complex manifold of dimension two with two algebraically
independent meromorphic functions is projective algebraic.
\source{hartshorne-algebraic-geometry:B.3.4}

\end{theorem}
\begin{proof}
The source states this theorem without proof and cites Chow and Kodaira.
\end{proof}

\section{Kähler Manifolds}
\label{ha-appB-sec4}
\source{hartshorne-algebraic-geometry:B.4}

\begin{definition}[Kähler, Hodge, and Moishezon manifolds]
\source{hartshorne-algebraic-geometry:B.4}
A Hermitian metric is Kähler when its associated $(1,1)$-form is closed.  A
compact Kähler manifold is a Hodge manifold if its Kähler class lies in
$H^2(X,\mathbb Z)$.  A compact complex manifold is Moishezon if its
meromorphic function field has transcendence degree equal to its dimension.
\end{definition}

\begin{theorem}[Kodaira]
\label{ha-appB-thm4-1}
\dcref{B.4.1}
Every Hodge manifold is projective algebraic.
\source{hartshorne-algebraic-geometry:B.4.1}

\end{theorem}
\begin{proof}
The source cites Kodaira for this theorem and does not give its proof.
\end{proof}

\begin{theorem}[Moishezon]
\label{ha-appB-thm4-2}
\dcref{B.4.2}
Every Moishezon manifold which is Kähler is projective algebraic.
\source{hartshorne-algebraic-geometry:B.4.2}

\end{theorem}
\begin{proof}
The source cites Moishezon for this theorem and does not give its proof.
\end{proof}

\section{The Exponential Sequence}
\label{ha-appB-sec5}
\source{hartshorne-algebraic-geometry:B.5}

\begin{construction}[Exponential sequence]
\source{hartshorne-algebraic-geometry:B.5}
The exponential sequence of sheaves on a complex manifold is
\[
0\longrightarrow\mathbb Z\longrightarrow\mathcal O_X
\xrightarrow{\exp(2\pi i\,\cdot)}\mathcal O_X^*\longrightarrow0.
\]
Its long exact sequence gives
\[
H^1(X,\mathcal O_X)\to \operatorname{Pic}X
\to H^2(X,\mathbb Z)\to H^2(X,\mathcal O_X),
\]
and $\operatorname{Pic}X\simeq H^1(X,\mathcal O_X^*)$.  For a curve of genus
$g$, $H^1(X,\mathcal O_X)\simeq\CC^g$ and
$\operatorname{Pic}^0(X)\simeq\CC^g/\mathbb Z^{2g}$, the Jacobian.
\end{construction}
